\chapter{Results}
\label{chap:results}

This compact results chapter keeps the quantitative evidence needed for the
final claim and moves the diagnostic gallery to Appendix~\ref{app:qualitative_diagnostics}.
The main reported model is the final prior-anchored residual GMM-head model,
evaluated on 2590 samples from 14 unseen test cities. All model-versus-prior
rows use the same held-out split and the same ground-only receiver mask.

\section{Milestones and prior quality}

Table~\ref{tab:timeline_results} condenses the experimental evolution. The
early rows are diagnostic, not directly comparable leaderboard entries, because
the mask contract and target formulation evolved during the project.

\begin{table}[htbp]
\centering
\caption{Representative experimental milestones retained in the compact TFG.}
\label{tab:timeline_results}
\scriptsize
\setlength{\tabcolsep}{3pt}
\begin{tabularx}{\textwidth}{p{0.24\textwidth}p{0.36\textwidth}X}
\toprule
Family & Lesson & Representative result \\
\midrule
Early U-Net/cGAN & Dense image-to-image prediction was feasible, but visual realism did not control physical RMSE. & Best clean pre-prior baseline around \SI{16.72}{dB}. \\
Ground-mask reset & Building pixels are not receivers and must be excluded. & Later metrics use the ground-only mask. \\
PMNet/PMHHNet & Height FiLM, physical channels and experts helped LoS, but NLoS tails remained poorly modelled. & PMHHNet reference: \SI{9.41}{dB} overall, \SI{3.75}{dB} LoS, \SI{31.48}{dB} NLoS. \\
Distribution-first GMM & The hard NLoS problem was distributional, not just architectural. & NLoS PL pixel-weighted RMSE around \SI{3.34}{dB} without explicit priors. \\
Frozen priors & Train-only calibrated priors became strong baselines in their own right. & \SI{1.9246}{dB} standalone PL prior (\SI{1.9383}{dB} final-test frozen PL); \SI{28.10}{ns} DS; \(13.74^\circ\) AS. \\
Final residual GMM head & A bounded neural residual improves already strong priors. & Test: \SI{1.6519}{dB} PL, \SI{26.5570}{ns} DS, \(11.3854^\circ\) AS. \\
\bottomrule
\end{tabularx}
\end{table}

The final path-loss prior is the most important baseline. Its LoS branch is not
FSPL alone: the enhanced two-ray calibration reduces LoS RMSE from about
\SI{3.89}{dB} to \SI{1.74}{dB}. The full hybrid prior reaches
\SI{1.9072}{dB} overall on the held-out validation+test split
(Table~\ref{tab:try78_results}). When re-evaluated on the final 14-city test
split as the frozen anchor for the neural model, the path-loss prior is
\SI{1.9383}{dB}.

\begin{table}[htbp]
\centering
\caption{Standalone calibrated path-loss prior on the held-out validation+test split.}
\label{tab:try78_results}
\scriptsize
\begin{tabular}{lcc}
\toprule
Region & RMSE & Unit \\
\midrule
LoS & 1.7516 & dB \\
NLoS & 3.3967 & dB \\
Overall & 1.9246 & dB \\
\bottomrule
\end{tabular}
\hfill
\begin{tabular}{lcc}
\toprule
LoS prior stage & LoS RMSE & Gain \\
\midrule
FSPL only & 3.9540 & N/A \\
FSPL + radial residual & 3.7086 & -0.25 \\
Enhanced two-ray & 1.7516 & -1.96 \\
\bottomrule
\end{tabular}
\end{table}

The spread priors are also strong enough to anchor the neural model. They are
fitted in \(\log(1+y)\) space and use the 114-key DS/AS regime dictionary
described in Chapter~3.

\begin{table}[htbp]
\centering
\caption{Calibrated spread-prior test results, pixel-weighted RMSE.}
\label{tab:try79_results}
\scriptsize
\begin{tabular}{lccc}
\toprule
Metric & Overall & LoS & NLoS \\
\midrule
Delay spread (ns) & 28.10 & 27.49 & 29.62 \\
Angular spread (\(^\circ\)) & 13.74 & 15.28 & 8.66 \\
\bottomrule
\end{tabular}
\end{table}

\section{Final unseen-city test results}

Table~\ref{tab:try80_final_test_overall} is the main thesis result. Negative
\(\Delta\) means that the final residual model improves the frozen prior. The
model improves RMSE for all three targets. Path loss is the cleanest success
case because both RMSE and MAE improve. Delay spread improves RMSE while MAE
increases slightly, consistent with reducing rare high-delay errors while
introducing small shifts in common low-delay regions. Angular spread improves
RMSE substantially and is essentially neutral in MAE.

\begin{table}[htbp]
\centering
\caption{Final model test results on 14 unseen cities.}
\label{tab:try80_final_test_overall}
\scriptsize
\setlength{\tabcolsep}{3pt}
\begin{tabular}{lrrrrrrl}
\toprule
Target & Model RMSE & Prior RMSE & \(\Delta\)RMSE & Model MAE & Prior MAE & \(\Delta\)MAE & Unit \\
\midrule
Path loss & 1.6519 & 1.9383 & -0.2865 & 0.9298 & 1.2319 & -0.3020 & dB \\
Delay spread & 26.5570 & 28.1023 & -1.5453 & 7.9897 & 7.3033 & +0.6864 & ns \\
Angular spread & 11.3854 & 13.7416 & -2.3562 & 5.0088 & 4.9890 & +0.0198 & deg \\
\bottomrule
\end{tabular}
\end{table}

The original project targets were \SI{5}{dB}, \SI{50}{ns} and
\(20^\circ\). The final model clears all three on unseen cities, while using
one joint network for PL, DS and AS.

\begin{table}[htbp]
\centering
\caption{Final test RMSE separated by geometric visibility.}
\label{tab:try80_final_test_los_nlos}
\scriptsize
\begin{tabular}{llrrr}
\toprule
Target & Scope & Model & Prior & \(\Delta\) \\
\midrule
PL & LoS & 1.4675 & 1.7519 & -0.2844 \\
PL & NLoS & 3.1482 & 3.5143 & -0.3661 \\
DS & LoS & 25.4419 & 27.4851 & -2.0432 \\
DS & NLoS & 29.2045 & 29.6157 & -0.4112 \\
AS & LoS & 12.3510 & 15.2818 & -2.9308 \\
AS & NLoS & 8.4424 & 8.6614 & -0.2191 \\
\bottomrule
\end{tabular}
\end{table}

The LoS/NLoS split matters because a global improvement could hide a failure in
the physically hard subset. Table~\ref{tab:try80_final_test_los_nlos} shows
that the residual correction improves both visibility regimes for all targets.

\begin{table}[htbp]
\centering
\caption{Validation versus test check for the selected final checkpoint.}
\label{tab:try80_val_test_sanity}
\scriptsize
\begin{tabular}{lrrrr}
\toprule
Split & Samples & PL RMSE & DS RMSE & AS RMSE \\
\midrule
Validation & 2750 & 1.6157 & 22.1061 & 11.3782 \\
Test & 2590 & 1.6519 & 26.5570 & 11.3854 \\
\bottomrule
\end{tabular}
\end{table}

The delay-spread test value is worse than validation mainly because the test
city mix contains harder high-delay cases; the frozen delay prior worsens in
the same direction, so this is a split composition effect rather than a change
in evaluator.

\section{Height, topology and split audits}

Height is not a cosmetic metadata field. It changes direct distance,
two-ray phase, elevation angle, blocker clearance and the selected calibration
regime. Table~\ref{tab:height_topology_audit} keeps the height audit compact,
and Table~\ref{tab:split_group_audit} keeps the residual audit across splits
and morphology/height groups.

\begin{table}[htbp]
\centering
\caption{Compact height audit. Values are prior-only RMSE: PL uses the
held-out Try~78 validation+test split, while DS/AS use the final 14-city
Try~79/Try~80 test split.}
\label{tab:height_topology_audit}
\scriptsize
\setlength{\tabcolsep}{4pt}
\begin{tabular}{lrrrr}
\toprule
Height bin & PL prior & DS prior & AS prior & PL samples \\
\midrule
12--50 m & 2.180 & 39.90 & 15.69 & 1211 \\
50--150 m & 1.936 & 27.17 & 14.19 & 2634 \\
150--300 m & 1.680 & 8.70 & 10.39 & 786 \\
300--500 m & 1.627 & 3.70 & 8.85 & 362 \\
\bottomrule
\end{tabular}
\end{table}

Low UAV heights are hardest, especially for delay spread; the highest bin has
much smaller spread values because the centred UAV sees over more local clutter.

\begin{table}[htbp]
\centering
\caption{Compact residual audit. Values are model-minus-prior RMSE deltas; negative is better.}
\label{tab:split_group_audit}
\scriptsize
\setlength{\tabcolsep}{4pt}
\begin{tabular}{lrrrr}
\toprule
Scope & Samples & PL & DS & AS \\
\midrule
Train & 10840 & -0.2895 & -2.0157 & -2.3578 \\
Validation & 2750 & -0.2758 & -1.3433 & -2.4396 \\
Test & 2590 & -0.2865 & -1.5453 & -2.3562 \\
All-city group range & 1250--2156 & -0.2363 to -0.3602 & -0.5428 to -3.1170 & -1.2874 to -2.8587 \\
\bottomrule
\end{tabular}
\end{table}

The final group audit is kept as a range rather than nine separate rows because
all entries have the same interpretation: the same residual model improves the
frozen priors across open, mixed and dense morphology families and low, mid and
high antenna bins. The detailed per-group rows are retained in
\texttt{results.tex} and the generated cluster outputs.

\section{Comparison with nearby SOA metrics}

Table~\ref{tab:soa_result_comparison} uses the same compact framing as the
paper. The table is not a formal leaderboard because datasets, transmitters,
split policies and target definitions differ. It is a scale check for the
path-loss number and a reminder that the PL/DS/AS joint task is less commonly
reported.

\begin{table}[htbp]
\centering
\caption{Compact comparison with nearby path-loss-side metrics.}
\label{tab:soa_result_comparison}
\scriptsize
\setlength{\tabcolsep}{3pt}
\begin{tabularx}{\textwidth}{p{0.25\textwidth}p{0.25\textwidth}X}
\toprule
Reference & PL-side metric & Reading / caveat \\
\midrule
RadioUNet \cite{radiounet2020} & \SI{1.6}{dB}--\SI{3.1}{dB} equivalent PL RMSE & Final PL is in this band; no UAV height or spread targets. \\
RadioGUNet \cite{radiogunet2025} & \SI{1.304}{dB} DPM; \SI{1.936}{dB} IRT & Final PL lies between those values in a different outdoor setting. \\
PMNet / ICASSP \cite{pmnet2023,icassp2023challenge} & \SI{1.38}{dB} converted score & Final PL is close but higher under a different benchmark contract; scale context only. \\
ReVeal \cite{reveal2025} & \SI{1.95}{dB} outdoor sparse-measurement RMSE & Final PL is in the same range; sparse mapping rather than dense generation. \\
RMTransformer / PathFinder \cite{rmtransformer2025,pathfinder2025} & \(\approx\SI{1.82}{dB}\) / \(\approx\SI{1.19}{dB}\) converted references & Useful scale checks under different split assumptions. \\
Gao et al.\ \cite{gao2026} & \SI{5.59}{dB} ICASSP~2023; \SI{12.19}{dB} ITU~2024 & Final PL is lower by a large margin, but their corridor-weighted outdoor PL task has different inputs/splits. \\
AIRMap \cite{airmap2025} & Path-gain error \(<\SI{4}{dB}\) against ray tracing & Same digital-twin motivation, different target and map size. \\
\textbf{Final model} & \textbf{\SI{1.6519}{dB} PL; \SI{26.5570}{ns} DS; \(11.3854^\circ\) AS} & \textbf{Strict CKM city holdout, centred variable-height UAV, dense \(513\times513\) PL/DS/AS.} \\
\bottomrule
\end{tabularx}
\end{table}

\section{Runtime}

The generator was benchmarked against the supplied MATLAB ray-tracing path on
Barcelona layouts at full \(513\times513\) resolution. Both rows use the same
laptop hardware: AMD Ryzen~5~5600H CPU and NVIDIA RTX~3050~Ti Laptop GPU. The
surrogate includes LoS/NLoS mask generation, frozen priors, final neural
prediction and default numeric export after the checkpoint has been loaded.

\begin{table}[htbp]
\centering
\caption{Runtime comparison for one full-resolution Barcelona CKM sample.}
\label{tab:inference_time_comparison}
\scriptsize
\begin{tabularx}{\textwidth}{lXr}
\toprule
Method & Timed path & Time \\
\midrule
MATLAB RT & Pixel receivers outside buildings; PL/DS/AS parsed from paths; CPU path & \SI{99.89}{s} \\
Final generator & Masks, priors, PL/DS/AS prediction, numeric export; CUDA mixed precision & \SI{0.88}{s} \\
\bottomrule
\end{tabularx}
\end{table}

This is a \(113.7\times\) speed-up under the measured setup. The result is not
only a neural-network forward-pass benchmark; it measures the usable generator
path, including priors and export.

\section{Limitations}
\label{sec:limitations}

The CKM labels are quantized, especially for spread targets. The model outputs
continuous values, but the ground truth contains staircase-like maps and
integer-valued spread levels. For path loss, a final MAE of
\SI{0.9298}{dB} is already close to the \SI{1}{dB} label step; the RMSE of
\SI{1.6519}{dB} remains higher because it emphasizes the remaining tails.
Thus \SI{0.93}{dB} MAE on a dataset quantized at roughly \SI{1}{dB} is a
strong result, and a cleaner or less quantized dataset would likely reveal more
of the continuous surrogate's resolution. This caveat does not invalidate the
reported metrics, but it limits how much meaning should be assigned to very
small residual changes near quantized boundaries.

The spread outputs are less mature than path loss. The final model improves
spread RMSE over the frozen priors, but delay-spread MAE increases slightly and
the remaining errors are concentrated in rare high-spread events. Future work
should therefore evaluate explicit tail metrics and probabilistic calibration,
not only RMSE and MAE.
